% Capstone 2: Pohlig--Hellman. Port of «Мини-курсовая - алгоритм Полига-Хеллмана».
% The performance numbers come from the generated requirements.tex (make bench).
\documentclass{ca-capstone}
\capstonerefs{alpertron-dilog,gmpy2,sagemath,hps}
\begin{document}
% Generated by `make bench CAPSTONE=pohlig-hellman` (taskgen.capstone report) from
% bench/results.json: reference time x slack on the reference machine.
% Do not edit: change bench/config.yaml and re-run the benchmark.
\newcommand{\ReqSlack}{10}
\newcommand{\ReqBitsMin}{50}
\newcommand{\ReqBitsMax}{100}
\newcommand{\ReqQBitsMax}{40}
\newcommand{\ReqSeconds}{20}

\capstonetitle{Решение задачи DLP алгоритмом Полига"--~Хеллмана}

\goal
Требуется реализовать метод решения задачи дискретного логарифмирования в
$\GF{p} \cong \Zn{p}$ алгоритмом Полига"--~Хеллмана.

\functionality
\begin{enumerate}
  \item Решение \emph{DLP}, в которой порядок основания "--- простое число, методом Шенкса или
    перебором\footnote{перебором имеет смысл решать относительно <<маленькие>> \emph{DLP} "--- для
    небольших модулей~$p$}
  \item Решение \emph{DLP}, в которой порядок основания "--- степень простого числа, методом
    сведения к собственным подгруппам
  \item Решение \emph{DLP}, в которой порядок основания "--- составное число произвольного вида,
    путём сведения к набору задач, в которых порядок основания "--- простое число или степень
    простого числа\footnote{собственно алгоритм Полига"--~Хеллмана; редукция заключается в
    формировании набора более простых задач, которые необходимо решить методами Шенкса и
    сведением к подгруппам, а полученные решения <<склеить>> по КТО}
\end{enumerate}

\requirements
Требования к реализации:
\begin{enumerate}
  \item корректность решения задачи \emph{DLP}
  \item приемлемая скорость работы (см. комментарий ниже)
\end{enumerate}

\textbf{Корректность работы} можно проверять, самостоятельно генерируя задачи, либо же
прибегая к помощи сторонних решателей, например, \cite{alpertron-dilog,sagemath}.

В данной работе важна не скорость, а корректность реализации и соответствие требованиям.
Однако тестовые примеры с числами размеров примерно \textbf{\ReqBitsMin--\ReqBitsMax~бит}
(простые делители порядка основания "--- не более \ReqQBitsMax~бит) должны выполняться за
приемлемое время "--- несколько десятков секунд (не более \ReqSeconds~с).

\details
Так как цель работы "--- решение \emph{DLP}, в своей реализации \textbf{можно пользоваться
библиотеками} для работы с большими целыми числами, а также пакетами линейной алгебры.

Рекомендуется использовать язык \emph{Python} и систему компьютерной алгебры
\emph{SageMath}~\cite{sagemath}.

Желательно сохранять промежуточные результаты решения \emph{DLP} алгоритмом Полига"--~Хеллмана,
такие как система сравнений, получаемая в результате выполнения редукции исходной задачи,
решения <<маленьких>> задач и проч.

Платформо-независимый, \textbf{переносимый код} писать \textbf{не требуется}, но во время сдачи
должна быть возможность проверить работу.

\schedule
Работа выполняется \textbf{индивидуально}. Срок выполнения "--- \textbf{месяц}.

Работа должна быть ваша: реализация должна быть \textbf{разработана и написана вами}, а не
собрана из кусков чужих работ.

Желательно хранить код в \emph{git}-репозитории на \emph{github}, \emph{bitbucket} или подобных
платформах. \textbf{Документация не требуется}.

\remarks
\begin{enumerate}
  \item Рекомендуется сначала реализовать метод решения \emph{DLP} Шенксом, затем собственно
    редукцию Полига"--~Хеллмана, и только после этого оптимизировать решение для случая, когда
    порядок показателя является степенью простого числа, реализуя метод сведения к
    собственным подгруппам
  \item Пусть это и крайне нежелательно, но, в случае нехватки времени, метод сведения к
    собственным подгруппам можно не реализовывать\footnote{но это может повлиять на итоговое
    <<ощущение от работы>>}
\end{enumerate}

\references
\end{document}
